% ============================================================
%  Limitations  (unnumbered; included after Conclusion via \section*)
% ============================================================
\label{sec:limitations}

\paragraph{Validation-set proxy.}
We promote candidates by training-subset accuracy with a guard against net regressions
(Eq.~\ref{eq:selection}), rather than by a held-out validation score. The guard ensures
monotonic training-set accuracy but does not rule out per-case regressions (a candidate
can fix some previously-incorrect cases while flipping others), and the selection-induced
bias from picking $c^{\star}$ over a discrete candidate set scales logarithmically in $M$.
We keep $M$ at a single-digit budget so this term stays small in practice, but quantifying
it empirically across benchmarks remains open.

\paragraph{Diagnosis bounded by the backbone.}
The \textsc{AnalyzerDecider}'s bottleneck classification (\emph{cognitive} vs.\
\emph{perceptual}) and the \textsc{Generator}'s candidates are produced by the same frozen
VLM that the agent uses. On backbones with weak self-introspection a mis-classified
bottleneck can propagate to a mis-targeted skill or tool; the promotion guard rejects
under-performing candidates but does not correct upstream diagnosis errors, so smaller or
less capable backbones may converge more slowly or to a less useful library.

\paragraph{Library governance.}
The capability library grows monotonically across iterations and across distribution-shift
phases. We merge skills covering the same problem class to prevent duplicate inflation
(Section~\ref{sec:method_skill}), but the library has no formal size bound; long-running
deployments may need additional pruning, versioning, or stale-capability eviction policies
that we do not investigate.

\paragraph{Benchmark scope.}
We evaluate four visual-reasoning benchmarks plus GTA and the VTool-R1 / DeepEyes splits.
Broader coverage (video QA, 3D spatial reasoning, medical imaging) is needed before strong
generality claims about the cognitive / perceptual decision rule.

\paragraph{Reliance on a per-case correctness signal.}
Both the offline evolution loop and the online adaptation policy
(Sections~\ref{sec:method_loop}--\ref{sec:method_adaptation}) require a per-case
correctness signal to drive diagnosis and the multi-candidate selection rule. Our
experiments use ground-truth labels from the benchmark splits as this signal, which
corresponds to deployment settings with human-in-the-loop review, an automated
quality-assurance pipeline, or an LLM-based post-hoc verifier. Deployments that have no
access to any of these channels---i.e.\ fully unsupervised production traffic---fall
outside the scope of the current method; extending \sysname{} to that setting would
require a confidence-based or self-consistency-based feedback substitute that we leave
to future work.
